% P10-4c-Nachtrag, 14.09.2026: E44-01 nach Artefakt und Urteilskategorie differenziert.
% P10-4c, 14.09.2026: Ledger-Herleitung und Nachweisanschlüsse präzisiert.
% P10-4b, 14.09.2026: Begriffs-/Beleganschlüsse und abgegrenzte Prüfaussagen präzisiert.
% M16, 11.09.2026: Leserführung und Absatzlogik; Anforderungen und Anker erhalten.
% B-34: E44-04 am Original geprüft: gewichtete Fehlerwerte, unterschiedliche Artefakte, gleiche Urteile bei variierenden Begründungen.
% ============================================================
% §4.4 Evidenzgebundene Wissensgewinnung — v07 (09.09.)
% v07: Deterministische Segmentierung der resultierenden Produktions-
% verarbeitung von modellgestützter Claim-Bildung getrennt; kein
% unbelegter historischer Wechsel der Segmentierung behauptet.
% v06 (01.09., Kollegen-Runde 4.4 — kritisch geprüft, alle Punkte
% berechtigt): ① Bilanz-Prüfungen auf „bereits gebildete Struktur-
% elemente" begrenzt (keine implizite Vollständigkeitsgarantie)
% ② „quellgestützt" → „quellengebunden erzeugt" (Struktur ≠ Treue;
% Treue = Kap. 7) ③ Produktionsverarbeitung und Evaluations-Referenz
% im Text getrennt; die damalige Zuschreibung der Produktions-
% segmentierung als modellgestützt wurde in v07 korrigiert.
% ④ Wende wissenschaftlich formuliert („Designrichtung verschob sich
% … zur Strukturierung der Quelle selbst") ⑤ „jede Einheit behandeln"
% → „explizit ausweisen, wie sie behandelt wurde"; Ersatz-Satz auf
% „Bilanz- und Strukturprüfungen" harmonisiert.
% v05 (01.09., Story-Auffrischung auf Autor-Wunsch, notes-verifiziert):
%   ① Zwei-Prüfrichtungen-Asymmetrie EXPLIZIT (Artefakt→Quelle durch
%     Zerlegung handhabbar/closed-world · Quelle→Artefakt offen/open-
%     world; B28 + Fazit Z. 1425 wörtlich belegt).
%   ② Die Wende als Erkenntnis-Satz: „Nicht der Prüfprozess, sondern
%     die Quelle selbst benötigte Struktur."
%   ③ Prüfsicht-Logik: quellgestützt VON ENTSTEHUNG AN + Pflicht zur
%     ausdrücklichen Behandlung jeder Einheit.
%   ④ Was den alten Review ERSETZTE: det. Struktur-/Abdeckungsprüfungen
%     (kein stiller Elementverlust) + menschliche Adjudikation.
%   Antwort auf Autor-Frage „einer von beiden nahezu gelöst — wodurch?":
%   die ARTEFAKT-Achse, durch Zerlegung + deterministisches Parsen
%   (span 0), nicht durch Modellwahl. Claims entstanden zuerst auf der
%   ARTEFAKT-Seite (B25/B26); Quellseiten-Struktur = Ledger.
% v04 (01.09., Konkretisierungs-Runde): M6 reale Claim-Ketten-Vignette
% (REQ-13, NEUE ID E44-06, nur Produktionsläufe) + E1-Ausstiegs-Satz
% („Evidence Ledger" benannt). v02/v03: Kollegen-Reviews (E44-04-
% Kausalgrenze, E44-05-Transfersatz, „erhalten könnte").
%
% Grundlage: Kapitel-4-Blaupause §9 (Absatzfolge 1–8, A3) · das
% FERTIGE, kollegengeprüfte Beispiel-Pivot (beispiel-pivot-kap4.md,
% v2 nach Review 23.08.) als Grundtext · genese-dossier P2+P3 / Z4–Z5
% (AUDITIERT 31.08.).
%
% Anpassungen gegenüber dem Beispiel-Grundtext:
%   ① Anforderung A3 als herausgehobener Satz ergänzt (Blaupause-
%     Wortlaut §9).
%   ② Ein S3-Satz zur explorativen Vorstudie ergänzt (Dossier-P3-
%     Formel: Richtung über zwei Domänen, KEINE Zahlen, "nicht
%     unabhängig verifiziert" deklariert; §5.2-S3-Regel) — die
%     Entscheidung fiel nicht blind, das darf gesagt werden.
%   ③ Übergabesatz auf Blaupause §9 Punkt 8 erweitert (Qualität,
%     Quellentreue, Abdeckung → Kapitel 7).
%   ④ "Abschnitt 2.3.1"-Rückverweis beibehalten (dort quellenbelegt:
%     semantische Vollständigkeitsbewertung gegen offene Quellen).
%
% Audit-Anker (für die Anhang-Tabelle; ALLE committet ✓):
%   Review-Grenze: thesis-evidence/D1-direct-vs-topic (DirectReview
%     modellabhängig+verrauscht; TopicCoverage reproduzierbar) +
%     source-claim-coverage-spike ("Wurzel-Ursache: freie Generierung
%     ohne Provenienz" — wörtlich im committeten README).
%   Evidence-First-Richtung: thesis-evidence/evidence-first-spike
%     (2 Domänen; A'-Vergleichsarm dokumentiert mit Commit+Run-ID).
%   (Privates Notes-Fazit phase02B:1425 "Precision lösbar, Coverage
%    nicht" = Schreibgrundlage, kein Anker.)
%
% Regeln: keine Mess-Verben, keine Spike-Zahlen (S3); Systemname
% seit E1-Entscheid 01.09. EINMAL als Ausstiegs-Satz benannt („wird in
% Kapitel 5 als Evidence Ledger konkretisiert"), Architektur bleibt
% Kap. 5; "nicht fortgeführter Review" ehrlich deklariert.
% Budget: ~1,4 Seiten (Ziel 1,3–1,6).
% ============================================================

\section{Evidenzgebundene Wissensgewinnung}
\label{sec:evidenzgebundene-wissensgewinnung}

\emph{Ausgangslage.} Der Strategievergleich benötigte ein
Prüfverfahren: Ein nachgelagerter Reviewprozess sollte die erzeugten
Artefakte gegen das Rohtranskript prüfen und feststellen, welche
Inhalte fehlten oder nicht durch die Quelle gedeckt waren. Zunächst
als Messinstrument gedacht, wurde dieser Prozess selbst zum
Untersuchungsgegenstand.

\emph{Vorgehen und formative Befunde.} Das Problem betraf sowohl
die erzeugten Inhalte als auch ihre Prüfung. In einem frühen
Kundenportal-Testfall ließ etwa das mit GPT-4.1-mini erzeugte
Architekturartefakt die in der Quelle angesprochenen Themen
Skalierbarkeit und Analytics, also die Auswertung von
Nutzungskennzahlen, unbehandelt. Diese beiden Lücken dienten später
als ausgewählte Prüffälle für den Reviewvergleich (E44-01).

Die freie Suche nach Fehlern und Auslassungen lieferte zudem
veränderliche Urteile. In einer weiteren Prüfanordnung stieg bei
zwei Bewertungen eines unveränderten Architekturartefakts der
gewichtete Fehlerwert von vier auf vierzehn; dies sind
Schweregradwerte, keine Befundzahlen (E44-04).
Ein daraufhin entwickelter Prototyp zerlegte ein Artefakt regelbasiert
in Prüfeinheiten, die ein Modell nach festen Kriterien gegen das
Transkript bewertete. Drei erhaltene Wiederholungen an einem
Anforderungsartefakt lieferten dieselben Einzelurteile bei
variierenden Begründungen. Da unterschiedliche Artefakte und
Prüfverfahren untersucht wurden, ist dies kein isolierter
Vorher-nachher-Nachweis (E44-04).

Dabei wurden zwei Prüfaufgaben erkennbar. \emph{Vom Artefakt zur
Quelle} war zu klären, ob eine erzeugte Aussage durch das Transkript
gedeckt war. Nach der Zerlegung standen die Aussagen fest und
konnten einzeln geprüft werden. \emph{Von der Quelle zum Artefakt}
lautete die Frage dagegen, welche relevanten Inhalte nicht übernommen
worden waren. Hier musste die Prüfung erst bestimmen, was enthalten
sein sollte. Ein Inhalt konnte sowohl bei der Erzeugung als auch
beim Review übersehen werden; ein ausbleibender Befund belegte
deshalb keine Vollständigkeit.

Eine feste Themenliste begrenzte diese zweite Prüfaufgabe. Im
Reviewvergleich wurden GPT-4.1-mini und GPT-5.4 als Prüfmodelle
eingesetzt. Im Architekturartefakt meldete die themengebundene
Prüfung mit GPT-5.4 beide ausgewählten Themen als fehlend.
GPT-4.1-mini meldete dort nur Skalierbarkeit als Lücke; Analytics
stufte es als teilweise abgedeckt ein, im Artefakt für offene Fragen
dagegen als fehlend. Die freie Prüfung mit GPT-5.4 übersah Analytics
in beiden Wiederholungen. Die gezielte Themenprüfung machte damit
beide Prüfgegenstände sichtbar, gewährleistete aber keine
übereinstimmende oder durchgehend richtige Einstufung (E44-01). Die Themenliste brachte jedoch
zusätzliche Information ein; eine allgemeine Modellunabhängigkeit
oder Überlegenheit folgte daraus nicht. Der Wechsel des Prüfmodells
ist zudem von einer erneuten Artefakterzeugung mit einem anderen
Modell zu unterscheiden.

Im damaligen Erzeugungsansatz wurde nicht systematisch mitgeführt,
welche Quellstellen welche Artefaktelemente begründeten. Ein
nachträglicher Reviewer musste daher rekonstruieren, was
berücksichtigt, abstrahiert oder ausgelassen worden war (E44-02).
Die Deutung als strukturelles Problem führte über die bloße
Verbesserung des Reviewers hinaus.

\emph{Entwurfsentscheidung.} Quellenbezug und überprüfbare
Zwischenstände sollten bereits während der Erzeugung entstehen.
Das Transkript sollte dafür zunächst in eine strukturierte,
quellgebundene Zwischenrepräsentation überführt werden. Explorative,
KI-gestützte Vorstudien über zwei Beispieldomänen erprobten die
Artefakterzeugung aus einem vorbereiteten, bestätigten Evidenzbestand.
Ihre Urteile deuteten auf eine bessere Erhaltung der Quellsemantik
hin und stützten die weitere Verfolgung dieses Ansatzes (E44-03).
Eingangsaufbereitung und Bewertungsverfahren unterschieden sich
jedoch zwischen den Varianten; die KI-Urteile wurden nicht unabhängig
verifiziert. Die Vorstudien begründen daher eine Designrichtung,
keinen isolierten Wirksamkeitsnachweis der späteren Ledger-Kette.

Die daraus entwickelte Verarbeitung segmentiert die Quelle
deterministisch in adressierbare Einheiten. Modelle bilden daraus
fachliche Aussagen, die \emph{Claims}, mit Belegstellen.
Deterministische Bilanz- und Strukturprüfungen kontrollieren,
welche Einheiten referenziert sind und ob bereits gebildete
Strukturelemente weitergegeben werden. Die Einordnung ungenutzter
Quellstellen bleibt eine semantische Aufgabe. Eine
\emph{Adjudikationsstufe} legt dem Menschen ausgewählte strittige
Evidenzzuordnungen zur Entscheidung vor. Damit werden mögliche
Auslassungen an ausdrücklich geführten Verarbeitungsergebnissen
untersuchbar; die vollständige Erfassung ihres Inhalts folgt daraus
nicht.

Nachgelagerte Artefaktelemente übernehmen die Claim-Referenzen.
Der Quellenbezug muss somit nicht erst aus dem fertigen Text
erschlossen werden. Ein dokumentierter Lauf des späteren Systems
illustriert dies: Aus der Interviewforderung nach einer
Profilübersicht mit anklickbaren Profilen entstand ein Claim mit
wörtlichem Quellzitat; das erzeugte Requirement übernahm seine
Kennung als Quellreferenz (E44-06). Der Fall zeigt eine
Rückverfolgungskette, keinen Qualitätsvergleich. Eine solche Referenz
ermöglicht Rückverfolgung, belegt aber noch nicht die inhaltliche
Stützung einer Aussage
(Abschnitt~\ref{sec:traceability-provenienz}).

Für die frühe Vollständigkeitsbewertung wurde außerdem eine von der
Systemverarbeitung getrennte Referenz vorbereitet. Die deterministische
Quellsegmentierung lieferte eine Annotationsvorlage; zunächst waren
elf Aussagen autorbestätigt. Eine umfangreichere Referenz mit
99 Einträgen wurde zurückgezogen, weil sie lediglich modellgeneriert
und nicht vom Autor bestätigt war (E44-05). Das spätere
Referenzverfahren beschreibt Abschnitt~\ref{sec:eval-gold}.

Der frühere Reviewprozess wurde im weiteren System nicht fortgeführt.
Sein Erkenntnisbeitrag lag darin, Inhaltsverluste und die Grenzen
ihrer nachträglichen Prüfung sichtbar zu machen. Die daraus
entwickelte Antwort verbindet semantische Modellarbeit mit expliziter
Quellenbindung, deterministischen Kontrollen und menschlicher
Klärung. Kapitel~\ref{chap:loesungsarchitektur} konkretisiert die
Zwischenrepräsentation als \emph{Evidence Ledger}. Daraus ergibt sich
die dritte Designanforderung:

\textbf{A3:} Fachliche Artefakte sollen aus einer expliziten, überprüfbaren
Evidenzstruktur erzeugt werden.

\emph{Einordnung und Konsequenz.} Begründet ist eine Antwort auf
die im untersuchten Entwicklungsstand erfahrenen Probleme, nicht
die allgemeine Notwendigkeit dieser Architektur. Die vollständige
Erhaltung relevanter Inhalte bleibt ein Ziel, an dem auch die
Zwischenrepräsentation zu prüfen ist. Kapitel~\ref{chap:evaluation}
bewertet deshalb Quellentreue und Inhaltsabdeckung des maschinellen
Claim-Bestands vor der menschlichen Adjudikation und untersucht die
gespeicherten Zwischenstände und Herkunftsbezüge. Die Inhaltsmessung
verwendet zwei bereits vorsegmentierte Quellenfälle; die Qualität
der vorgelagerten Zerlegung ist nicht Bestandteil dieser Messung.

% Historische Inline-Herkunftsnotizen vor M16; aktueller Befund im M16-Bericht.
% [Fundstelle Messinstrument-Motiv: fruehphase/qualitaetsrubrik.md Kopf —
%  „Zweck: Fehlerbasierte, transkript-gebundene Bewertung … für den
%  A/B/C-Vergleich"; Autor-Klärung 31.08. (Kausalkette)]
% [Fundstelle E44-01: phase02-notes Iter 23–26 (Jury-Kalibrierung +
%  Instrument-Analyse); phase02B-notes B4/B5 (Judge instabil), B19
%  (Head-to-Head „stark modellabhängig"); thesis-evidence/
%  D1-direct-vs-topic/evidence.md]
% [Fundstelle E44-04: phase02B-notes — B25 (Per-Item-Prototyp: span 0
%  statt ±10; 4/4/4 + 51/51 identisch; Run 20260617_162205_950a9a/
%  jury/_peritem; deterministisches Parsen „KEIN LLM → null Varianz"),
%  B26 (classify-units integriert), B27 (2. Transkript, nicht overfit)]
% [Fundstellen Achsen-Asymmetrie/Open-World: phase02B-notes Z. ~1425
%  (Fazit: „Precision (closed-world) lösbar; Recall/Coverage (open-world)
%  ist das Problem"), B28 (Coverage-Achse „rauer als die Artefakt-Achse"),
%  Commited-Notes B34 (Coverage als Score unbrauchbar) / B35 (Topic-
%  Aggregation 58→4); evaluation-uebergang/Bewertung-Evidence-Ledger.md §1;
%  Zwei-Achsen-Design wörtlich: B28 „Artefakt-Achse = solide/einsatzreif
%  (span 0) · Coverage-Achse = rauer erster Schnitt", Fazit Z. 1425
%  „Precision (closed-world, verifizierbar) ist lösbar; Recall/Coverage
%  (open-world, Urteil) ist das härtere Problem"; Commited-Notes Z. 93
%  „bestätigt das Zwei-Achsen-Design"]
% [Fundstelle E44-02: thesis-evidence/source-claim-coverage-spike/
%  evidence.md + thesis-evidence/README Spike-4-Zeile („Wurzel-Ursache
%  benannt: freie Generierung ohne Provenienz")]
% [Präzisiert 01.09. (Kollegen-Punkt 1, berechtigt): die Prüfungen
%  sichern die WEITERFÜHRUNG bereits gebildeter Elemente (Unit-Bilanz,
%  CANDIDATE_SILENTLY_DROPPED), NICHT die Transkript-Vollständigkeit —
%  die bleibt lt. Einordnung zu evaluieren.]
% [Kollegen-Punkte 2+5 (berechtigt): „quellengebunden" = Struktur-
%  Eigenschaft (Referenz vorhanden) ≠ „quellgestützt" (Quelle trägt
%  inhaltlich — erst Kap. 7); „ausweisen wie behandelt" statt
%  „behandeln" — auch Rauschen wird bewusst klassifiziert, nicht jede
%  Unit wird Artefaktinhalt.]
% [09.09. am Code geprüft: AtomicUnitSegmentationExecutor ruft
%  AtomicUnitSegmenter.Segment auf; dieser verwendet deterministisch
%  TranscriptSegmenter.Segment und vergibt positionale AU-Kennungen.
%  Die frühere LLM-Zuschreibung aus der README-Tabelle war hierfür
%  unzutreffend. Der Text beschreibt die resultierende Verarbeitung,
%  keinen nachträglich angenommenen historischen Methodenwechsel.]
% [Fundstelle E44-06 (M6, verifiziert 01.09., NUR Produktionsläufe):
%  Transkript input/transcripts/Interview-Einrichtung.txt Z. 223 →
%  Atomic Unit AU-0045 → Claim canon-profiluebersicht-nach-login mit
%  evidence.quote (runs/ledger/20260723_093443_a85a80/step-03b-
%  adjudicated/consumable.json) → REQ-13 sourceClaimIds
%  (runs/recipe/20260723_113115_35335a + state/core/project-state.json);
%  bewusst NICHT input/eval-labels (Referenz-/Produktions-Trennung)]
% [Produktionsverarbeitung und unabhängige Evaluationsreferenz bleiben
%  nach Zweck und Bewertungsverantwortung getrennt; deterministische
%  Segmentierung ist nicht auf die Referenzbildung beschränkt.]
% [Fundstelle E44-05: evaluation-uebergang/Bewertung-Evidence-Ledger.md
%  §2.2 — R1 („kein LLM ist Ground Truth"), R2 (erste nicht-zirkuläre
%  Zahl gegen autor-bestätigte 11er-Fixture), R3 (Rückzug der modell-
%  erzeugten 99er-Referenz), R4→L2 (Ausreißer-Korrektur); Ledger-Strang:
%  thesis-story/ledger-aera/; Fixture committet: input/eval-labels/]
% [Fundstelle E44-03: thesis-evidence/evidence-first-spike + README
%  Spike-5-Zeile (2 Domänen; A'-Arm Commit c34b9a0 /
%  Run 20260630_145558_68149f; Grenzen dort deklariert)]
% [Ersatz-des-Reviews verifiziert 01.09.: evaluation-uebergang/Bewertung-
%  Evidence-Ledger.md Z. 42 — Qualitäts-Gate prüft OHNE LLM (Schema/Enum/
%  Cluster-Invariante, Fehlschlag statt Durchreichen), „stiller
%  Kandidatenverlust von Warnung zu Fehler hochgestuft", CoverageRepair-
%  Executor, erster vollständig grüner Lauf (0/0) — Recall-PROZENT dort
%  bewusst NICHT zitiert (Autor-Regel). Design-Beschreibung des
%  Folgesystems, KEIN eigener formativer Anker (Nachweis = Kap. 5/6).]
% [E1-Entscheid 01.09.: resultierendes Konzept BENENNEN, Architektur
%  nicht erklären — Ausstiegs-Satz-Form]
